\begin{frame}
\frametitle{Hardwarová realizace registru}
Základní realizace registru pomocí tzv. obvodu d-latch
\begin{center}
\includegraphics[width=0.5\textwidth]{generated/cpu/register/d_latch.pdf}
\end{center}

Tento obvod je velkou změnou oproti dosavadnímu přístupu, protože se zde objevuje cyklus -- výstup hradla je vstupem stejného hradla, nebo vstupem hradla jehož výstup je vstupem prvního hradla.

Co to tedy dělá?

\end{frame}

\begin{frame}
\frametitle{Hardwarová realizace registru}

Chování obvodu je určeno hodnotou vstupu E.

Nejdříve se podíváme, co se děje při E=1:
\begin{columns}
\begin{column}{0.5\textwidth}
\begin{center}
\includegraphics[width=\textwidth]{generated/cpu/register/d_latch_e1-cz.pdf}
\end{center}
\end{column}
\begin{column}{0.5\textwidth}  
\begin{itemize}
\item Vstup D se se zpožděním objeví na výstupu Q
\begin{itemize}
\item pro $D=1$ musí být výstup $Q=1$, protože jeden vstup do posledního nand hradla je 0, tedy výstup musí být 1
\item pro $D=0$ musí být $\overline{Q}=1$, protože opět jeden vstup do nand hradla je 0, když je $\overline{Q}=1$, tak potom je výstup $Q=0$, protože oba dva vstupy jsou 1, výstup nand je 0
\end{itemize}
\end{itemize}
\end{column}
\end{columns}

\end{frame}

\begin{frame}
\frametitle{Hardwarová realizace registru}

Pokud je E=0:
\begin{columns}
\begin{column}{0.5\textwidth}
\begin{center}
\includegraphics[width=\textwidth]{generated/cpu/register/d_latch_e0-cz.pdf}
\end{center}
\end{column}
\begin{column}{0.5\textwidth}  
\begin{itemize}
\item Hodnota $Q,\overline{Q}$ závisí pouze na staré hodnotě dvojice $Q,\overline{Q}$
\item Dvojice nabývá pouze hodnot, která mohla vzniknout při zápisu, tedy buď $Q=1,\overline{Q}=0$ nebo $Q=0,\overline{Q}=1$
\end{itemize}
\end{column}
\end{columns}

\end{frame}

\begin{frame}
\frametitle{Hardwarová realizace registru}

Aktivace náběžnou hranou:
\begin{columns}
\begin{column}{0.45\textwidth}
\begin{center}
\includegraphics[width=0.9\textwidth]{generated/cpu/register/edge.pdf}
\end{center}
\end{column}
\begin{column}{0.55\textwidth}  
\begin{itemize}
\item Může být výstup tohoto obvodu vůbec někdy v 1?
\item matematicky je vždy $(C$~and~$\overline{C})=0$ 
\item ve skutečnosti je ale hradlo not zpožděno oproti normálnímu signálu a na tuto dobu zpoždění je výstup \texttt{E'} v 1
\item pokud by tato hodnota byla moc krátká na to, aby se obvod d-latch nastavil, tak je možné přidat negací více za sebou (zelená křivka)
\item tomuto obvodu se říká d flip-flop
\end{itemize}
\end{column}
\end{columns}

\end{frame}
